%-------------------------------------------------------------------------------
% 3. 2026 DACON — 스마트 창고 출고 지연 예측 AI 경진대회 (상위 11%)
%-------------------------------------------------------------------------------
\cvsection{3. 2026 DACON : 스마트 창고 출고 지연 예측 AI 경진대회}

\projsummary{AMR(자율이동로봇) 스마트 물류창고 운영 데이터를 기반으로, 출고 지연 시간을 미리 예측하는 모델 개발}
\projfig{p3.png}

\projpart{배경 및 목표}
\begin{cvitems}
  \item {\textbf{배경}\enskip 로봇으로 운영되는 스마트 창고의 출고 지연은 로봇·주문량·혼잡도 등 다수 요인이 얽혀 발생 $\to$ 사전 예측으로 선제 조치가 가능한 모델이 필요.}
  \item {\textbf{목표}\enskip 12,000개의 독립적인 창고 운영 시나리오를 학습해 \textbf{향후 30분간의 평균 출고 지연 시간(분)}을 예측 (평가 지표: MAE).}
  \item {\textbf{개발 기간}\enskip 2026.04.01 -- 2026.05.04}
  \item {\textbf{기술 스택}\enskip Python · PyTorch · LightGBM · CatBoost · XGBoost · Bi-GRU · MultiheadAttention · Optuna · SciPy(SLSQP) · Dash}
\end{cvitems}

\projpart{데이터 소개}
\begin{cvitems}
  \item {시나리오 12,000개 $\times$ 25 timestep 운영 로그 + 창고 300개 정적 메타정보(통로 폭·교차로 수·포장대 등).}
  \item {\textbf{분포 특성}\enskip Train 250 / Test 100 / 교집합 50 / Test에만 50 $\to$ \textbf{신규 창고 일반화}가 핵심.}
\end{cvitems}

\projpart{설계 과정 및 수행 내용}
\begin{cvitems}
  \item {\textbf{\cn{1} 시계열 구조 복원 + AMR 도메인 피처 설계}\enskip 시나리오 안 행 순서가 뒤섞여 시계열 학습이 불가하여 시나리오 ID·행 번호로 재정렬해 시간 순서를 복원하고 시점 인덱스(0--24)를 피처화. 핵심 운영 지표 8개에 시간 변형 4종(lag·rolling·pure-past·시나리오 평균)을 적용하고, AMR 도메인 합성 피처(교차로 지연 비용, 단계별 부하 지표, 실질 혼잡도 = 혼잡도 $\times$ 교차로 밀도 / 통로 폭)로 '좁은 통로일수록 같은 혼잡도가 더 큰 지연을 만든다'는 도메인 지식을 수식화.}
  \item {\textbf{\cn{2} 신규 창고 일반화: kNN 인코딩 + GroupKFold}\enskip 학습에 없던 창고 50개에서도 잘 예측하도록 두 장치를 병행. \textbf{kNN 인코딩} — 창고 정적 정보 13개(통로 폭·교차로 수·포장대 수·면적 등)와 구조 타입을 정규화해 코사인 거리로 유사한 5개 창고의 평균·표준편차·최댓값을 피처화(자기 자신 제외로 유출 차단). \textbf{GroupKFold(by layout\_id)} — 같은 창고가 학습·검증에 동시 등장하지 않게 분할해 신규 창고만으로 검증.}
  \item {\textbf{\cn{3} 다중 모델 앙상블 + SLSQP 가중치 최적화}\enskip 트리 모델만으로는 25개 시간 구간의 흐름을 학습하기 어려워 트리 3종과 시퀀스 2종을 함께 구성(XGBoost는 \cn{4} 진단 후 제외, 최종 4모델). 트리 3종(LightGBM·CatBoost·XGBoost)은 Optuna로 GroupKFold 하이퍼파라미터를 탐색하고 비대칭 분포 보정을 위해 log1p 학습·expm1 복원. Bi-GRU는 시나리오를 (25, N) 시퀀스로 재구성해 LSTM 대비 경량으로 일반화에 유리(LB 10.1669 $\to$ 10.0194), Bi-GRU + Multi-Head Attention은 4-head·잔차 연결·LayerNorm으로 멀리 떨어진 시점 관계를 포착. SLSQP로 OOF 최소화 가중치를 자동 탐색(합=1, 0--1).}
  \item {\textbf{\cn{4} OOF--LB 디커플링 진단 및 모델 재구성}\enskip GroupKFold 이후에도 OOF--LB 갭 1.4가 지속되고 'OOF 개선 / LB 악화'가 4--5회 반복되어 검증 지표 자체의 구조적 문제로 판단. \textbf{(1) Adversarial Feature Selection} — train/test 이진 분류기가 AUC = 1.0으로 완전 분리됨을 확인하고 피처별 누수 점수를 정량화해 5개 제거(트리 3종만 재학습, GRU/Attention은 시퀀스 신호 보존을 위해 원본 유지) $\to$ LB 9.9983 $\to$ 9.9917. \textbf{(2) 앙상블 가중치 기반 노이즈 모델 식별} — SLSQP에서 XGBoost 가중치가 0.0100으로 수렴해 노이즈로 판단·제외하고 4모델로 재구성(LightGBM 0.20 $\to$ 0.28, CatBoost 0.16 $\to$ 0.17, GRU 0.29, Attention GRU 0.26) $\to$ LB 9.9917 $\to$ 9.9890.}
  \item {\textbf{\cn{5} Dash 기반 운영 대시보드 구현}\enskip '예측 $\to$ 진단 $\to$ 시뮬레이션 $\to$ 액션'을 한 화면에 잇는 6개 섹션(물류 파이프라인 $\to$ 현재 예측 $\to$ 시점별 추이·알림 $\to$ SHAP 기여 분해 $\to$ 운영 패턴 Radar $\to$ What-If 시뮬레이션)을 설계. What-If 슬라이더로 핵심 신호 4개를 조정하면 연동된 16개 합성 피처가 자동 재계산되어 즉시 새 예측을 표시. 학습 노트북의 피처 엔지니어링을 모듈화하고 50,000행 캐시 + 20개 모델(5-fold $\times$ 4종) + SHAP 사전 계산(50,000 $\times$ 168)으로 '30분 후 위험 $\to$ 원인 신호 $\to$ 액션 효과'를 한 페이지에서 검증.}
\end{cvitems}

\projpart{성과}
\begin{cvparagraph}
베이스라인 MAE \textbf{12.6160 $\to$ 9.9890} 개선, \textbf{최종 68위 (상위 11\%)}. 검증 지표의 구조적 결함을 진단·교정한 과정이 성능 향상의 핵심.
\end{cvparagraph}
